\begin{definition*}
Let $X$ be a set. A metric on $X$ is a function
\[
  d : X \times X \to \mathbb{R}
\]
such that, for all $x,y,z \in X$,
\[
  d(x,y) \ge 0,
\]
\[
  d(x,y) = 0 \iff x = y,
\]
\[
  d(x,y) = d(y,x),
\]
and
\[
  d(x,z) \le d(x,y) + d(y,z).
\]
These four requirements are called the metric axioms: nonnegativity, identity
of indiscernibles, symmetry, and the triangle inequality.

In Lean, we formalize this definition as a structure:
\[
\begin{gathered}
  \texttt{structure Metric (Point : Type u) where}\\
  \texttt{\ \ distance : Point -> Point -> Real}\\
  \texttt{\ \ distance\_nonnegative : forall x y : Point, 0 <= distance x y}\\
  \texttt{\ \ distance\_eq\_zero\_iff\_equal :}\\
  \texttt{\ \ \ \ forall x y : Point, distance x y = 0 <-> x = y}\\
  \texttt{\ \ distance\_symmetric : forall x y : Point, distance x y = distance y x}\\
  \texttt{\ \ distance\_triangle :}\\
  \texttt{\ \ \ \ forall x y z : Point, distance x z <= distance x y + distance y z.}
\end{gathered}
\]
The field \texttt{distance} is the proposed distance function. The remaining
fields are the proof obligations saying that this function satisfies the metric
axioms. Thus, constructing a value of type \texttt{Metric Point} means providing
both a formula for distance and proofs that the formula behaves like a metric.
\end{definition*}

\begin{example*}
The real line $\mathbb{R}$ carries the usual absolute-value metric
\[
  d(x,y) = |x-y|.
\]
In Lean, this example begins as
\[
\begin{gathered}
  \texttt{def realAbsoluteValueMetric : Metric Real where}\\
  \texttt{\ \ distance x y := |x - y|}\\
  \texttt{\ \ distance\_nonnegative := by sorry}\\
  \texttt{\ \ distance\_eq\_zero\_iff\_equal := by sorry}\\
  \texttt{\ \ distance\_symmetric := by sorry}\\
  \texttt{\ \ distance\_triangle := by sorry.}
\end{gathered}
\]
The distance formula is now fixed. What remains is to prove that this formula
satisfies each metric axiom. In other words, for arbitrary real numbers
$x,y,z$, we must show:
\[
  0 \le |x-y|,
\]
\[
  |x-y| = 0 \iff x = y,
\]
\[
  |x-y| = |y-x|,
\]
and
\[
  |x-z| \le |x-y| + |y-z|.
\]
These are the next proof tasks for the real-line metric.

For nonnegativity, Mathlib contains the theorem \texttt{abs\_nonneg}. Its
mathematical content is that the absolute value of any real number is
nonnegative:
\[
  0 \le |a|.
\]
The metric axiom for the real-line distance asks us to prove
\[
  0 \le |x-y|
\]
for arbitrary real numbers \texttt{x} and \texttt{y}. The Lean proof is:
\[
\begin{gathered}
  \texttt{distance\_nonnegative := by}\\
  \texttt{\ \ intro x y}\\
  \texttt{\ \ exact abs\_nonneg (x - y).}
\end{gathered}
\]
The line \texttt{intro x y} introduces arbitrary real numbers \texttt{x} and
\texttt{y}. After that, the goal is
\[
  \texttt{0 <= |x - y|}.
\]
This is exactly \texttt{abs\_nonneg} applied to the real number
\texttt{x - y}. So \texttt{exact abs\_nonneg (x - y)} closes the goal.

For identity of indiscernibles, the metric axiom asks us to prove
\[
  |x-y| = 0 \iff x = y.
\]
Because the goal is an if-and-only-if statement, the proof naturally splits
into two directions. In Lean, the command \texttt{constructor} turns a goal of
the form \texttt{A <-> B} into two goals: first \texttt{A -> B}, then
\texttt{B -> A}. The proof begins:
\[
\begin{gathered}
  \texttt{distance\_eq\_zero\_iff\_equal := by}\\
  \texttt{\ \ intro x y}\\
  \texttt{\ \ constructor.}
\end{gathered}
\]
The line \texttt{intro x y} introduces arbitrary real numbers \texttt{x} and
\texttt{y}. Then \texttt{constructor} asks us to prove both implications in the
iff.

For the forward direction, the hypothesis is
\[
  \texttt{hypothesis : |x - y| = 0}
\]
and the goal is \texttt{x = y}. Mathlib's theorem \texttt{abs\_eq\_zero} says
that an absolute value is zero exactly when the expression inside it is zero.
Using \texttt{.mp}, the forward direction of that iff, gives
\[
  \texttt{x - y = 0}.
\]
From there, \texttt{linarith} proves the resulting linear arithmetic goal
\texttt{x = y}. In Lean:
\[
\begin{gathered}
  \texttt{-- =>}\\
  \texttt{\ \cdot\ intro hypothesis}\\
  \texttt{\ \ \ have difference\_eq\_zero : x - y = 0 :=}\\
  \texttt{\ \ \ \ abs\_eq\_zero.mp hypothesis}\\
  \texttt{\ \ \ linarith.}
\end{gathered}
\]
The line beginning with \texttt{--} is a one-line Lean comment. It is ignored by
Lean and is only there to label the forward direction of the proof.

For the reverse direction, the hypothesis is
\[
  \texttt{hypothesis : x = y}
\]
and the goal is \texttt{|x - y| = 0}. Rewriting by \texttt{hypothesis} changes
\texttt{x - y} into \texttt{y - y}; then \texttt{sub\_self} simplifies this to
\texttt{0}. The remaining goal is \texttt{|0| = 0}, which is exactly
\texttt{abs\_zero}. In Lean:
\[
\begin{gathered}
  \texttt{-- <=}\\
  \texttt{\ \cdot\ intro hypothesis}\\
  \texttt{\ \ \ rw [hypothesis, sub\_self]}\\
  \texttt{\ \ \ exact abs\_zero.}
\end{gathered}
\]
Here the comment \texttt{-- <=} labels the reverse direction of the iff.

For symmetry, Mathlib contains the theorem \texttt{abs\_sub\_comm}. Its
mathematical content is
\[
  |x-y| = |y-x|.
\]
This is exactly the symmetry axiom for the absolute-value distance on
\(\mathbb{R}\). The Lean proof is:
\[
\begin{gathered}
  \texttt{distance\_symmetric := by}\\
  \texttt{\ \ intro x y}\\
  \texttt{\ \ exact abs\_sub\_comm x y.}
\end{gathered}
\]
The line \texttt{intro x y} introduces arbitrary real numbers \texttt{x} and
\texttt{y}. After that, the goal is
\[
  \texttt{|x - y| = |y - x|},
\]
which is precisely what \texttt{abs\_sub\_comm x y} proves.

For the triangle inequality, Mathlib already contains the theorem
\texttt{abs\_add\_le}, whose mathematical content is
\[
  |a+b| \le |a| + |b|.
\]
The documentation describes this as saying that absolute value satisfies the
triangle inequality. However, Lean can only use that theorem when the current
goal has been put into a matching form. It is not enough for us to know
informally that absolute value has the triangle inequality; the expression in
the proof state must be connected to the theorem's expression.

For the real-line metric, the triangle axiom asks us to prove
\[
  |x-z| \le |x-y| + |y-z|.
\]
This does not syntactically look like \texttt{|a + b| <= |a| + |b|}. The
algebraic bridge is
\[
  x-z = (x-y) + (y-z).
\]
After rewriting the left-hand side using this identity, the goal becomes
\[
  |(x-y)+(y-z)| \le |x-y| + |y-z|,
\]
which is exactly the theorem \texttt{abs\_add\_le} applied to
\[
  a = x-y,
  \qquad
  b = y-z.
\]
Thus the proof has two conceptual steps: first rewrite the distance
\texttt{|x - z|} into an absolute value of a sum, and then apply
\texttt{abs\_add\_le}. This is a common pattern in Lean formalization: a
library theorem may express the right mathematical fact, but we still have to
arrange the goal so that Lean can see the theorem applies.

The corresponding Lean proof can be written as follows:
\[
\begin{gathered}
  \texttt{distance\_triangle := by}\\
  \texttt{\ \ intro x y z}\\
  \texttt{\ \ calc}\\
  \texttt{\ \ \ \ |x - z| = |(x - y) + (y - z)| := by}\\
  \texttt{\ \ \ \ \ \ ring\_nf}\\
  \texttt{\ \ \ \ \_ <= |x - y| + |y - z| := by}\\
  \texttt{\ \ \ \ \ \ exact abs\_add\_le (x - y) (y - z).}
\end{gathered}
\]
The line \texttt{distance\_triangle := by} begins the proof of the
\texttt{distance\_triangle} field of the \texttt{Metric} structure. Since the
distance on the real line has been defined by \texttt{distance x y := |x - y|},
the expected proof is
\[
  \texttt{forall x y z : Real, |x - z| <= |x - y| + |y - z|}.
\]
The command \texttt{intro x y z} introduces arbitrary real numbers
\texttt{x}, \texttt{y}, and \texttt{z}. This is the Lean version of saying
``Let $x,y,z \in \mathbb{R}$ be arbitrary.'' After these introductions, the
goal is
\[
  \texttt{|x - z| <= |x - y| + |y - z|}.
\]

The command \texttt{calc} starts a calculation chain. We prove the desired
inequality by moving from the left-hand side to the right-hand side through an
intermediate expression:
\[
  |x-z|
  =
  |(x-y)+(y-z)|
  \le
  |x-y|+|y-z|.
\]
The first step,
\[
  \texttt{|x - z| = |(x - y) + (y - z)| := by ring\_nf,}
\]
is just algebra. The tactic \texttt{ring\_nf} normalizes ring expressions and
proves identities such as
\[
  x-z = (x-y)+(y-z).
\]
To normalize a ring expression means to rewrite it into a standard, canonical
algebraic form: like terms are collected, parentheses are removed according to
associativity and distributivity, subtractions are treated as additions of
negatives, and coefficients are simplified. For example, both
\[
  x-z
  \qquad\text{and}\qquad
  (x-y)+(y-z)
\]
normalize to the same formal expression, namely \(x-z\). Once the two sides
have the same ring normal form, Lean can recognize them as equal. Because the
expressions inside the absolute values are algebraically equal, Lean can use
this normalization to prove the equality of the absolute values.

The second step uses an underscore:
\[
  \texttt{\_ <= |x - y| + |y - z| := by exact abs\_add\_le (x - y) (y - z).}
\]
In a \texttt{calc} block, the underscore means ``the expression from the
previous line.'' Thus this line is proving
\[
  \texttt{|(x - y) + (y - z)| <= |x - y| + |y - z|}.
\]
That is exactly the statement of \texttt{abs\_add\_le} with
\[
  a = x-y,
  \qquad
  b = y-z.
\]
The command \texttt{exact abs\_add\_le (x - y) (y - z)} tells Lean that this
library theorem proves the current goal exactly. The whole calculation gives
the metric triangle inequality for the absolute-value distance on
\(\mathbb{R}\).

At this point, the real-line metric proof in Lean is:

\begin{verbatim}
/-- The usual absolute-value metric on the real line. -/
def realAbsoluteValueMetric : Metric Real where
  distance x y := |x - y|

  distance_nonnegative := by
    intro x y
    exact abs_nonneg (x - y)

  distance_eq_zero_iff_equal := by
    intro x y
    constructor
    -- =>
    · intro hypothesis
      have difference_eq_zero : x - y = 0 := abs_eq_zero.mp hypothesis
      linarith
    -- <=
    · intro hypothesis
      rw [hypothesis, sub_self]
      exact abs_zero

  distance_symmetric := by
    intro x y
    exact abs_sub_comm x y

  distance_triangle := by
    intro x y z
    calc
      |x - z| = |(x - y) + (y - z)| := by
        ring_nf
      _ ≤ |x - y| + |y - z| := by
        exact abs_add_le (x - y) (y - z)
\end{verbatim}
\end{example*}

\begin{example*}
Every set also carries the discrete metric. Mathematically, this metric is
defined by
\[
  d(x,y)
  =
  \begin{cases}
    0, & x = y,\\
    1, & x \ne y.
  \end{cases}
\]
The idea is that distinct points are all placed exactly one unit apart.

In Lean, the definition uses an \texttt{if}-expression:
\[
  \texttt{distance x y := if x = y then 0 else 1.}
\]
To write this for an arbitrary type \texttt{Point}, Lean needs the assumption
\[
  \texttt{[DecidableEq Point]}.
\]
This says that equality of points is decidable: for any two points, Lean can
choose the \texttt{then} branch or the \texttt{else} branch of the
\texttt{if}-expression. This is the formal version of being able to determine
whether \(x=y\) or \(x\ne y\).

The Lean definitions are:

\begin{verbatim}
/-- The discrete metric on any type of points with decidable equality. -/
def discreteMetric (Point : Type u) [DecidableEq Point] : Metric Point where
  distance x y := if x = y then 0 else 1
  distance_nonnegative := by
    sorry
  distance_eq_zero_iff_equal := by
    sorry
  distance_symmetric := by
    sorry
  distance_triangle := by
    sorry

/-- Any type with decidable equality equipped with the discrete metric. -/
def discreteMetricSpace (Point : Type u) [DecidableEq Point] : MetricSpace Point where
  metric := discreteMetric Point
\end{verbatim}

The proof obligations are left as exercises. They will mostly involve case
splits. The command \texttt{by\_cases} introduces one branch where an equality
holds and another branch where it does not. In each branch, \texttt{simp} can
reduce the \texttt{if}-expression to either \texttt{0} or \texttt{1}. For
example, if \texttt{x = y}, then \texttt{if x = y then 0 else 1} simplifies to
\texttt{0}; if \texttt{x ≠ y}, it simplifies to \texttt{1}. The triangle
inequality is longer only because there are three relevant equalities:
\texttt{x = z}, \texttt{x = y}, and \texttt{y = z}.
\end{example*}
